\documentclass[blue]{beamer}

\mode<presentation> {
  \setbeamertemplate{background canvas}[vertical shading][bottom=red!10,top=blue!10]

  \usetheme{Boadilla}
  \usefonttheme[onlysmall]{serif}
}


\usepackage{pgf,pgfarrows,pgfnodes,pgfautomata,pgfheaps,pgfshade}
\usepackage{amsmath,amssymb,bm}
\usepackage{colortbl}
\usepackage[english]{babel}
\usepackage[latin1]{inputenc}

\setbeamercovered{dynamic}

\DeclareMathOperator\PV{PV}

\def\R{\mathbb{R}}
\def\Q{\mathbb{Q}}
\def\N{\mathbb{N}}
\def\E{\PV}
\newcommand{\ps}[2]{ \langle #1 , #2 \rangle }
\def\1{\mathbf{1}}
\newcommand{\nor}[1]{\left\| #1 \right\|}
\def\leq{\leqslant}
\def\geq{\geqslant}
\def\et{\quad \textrm{and} \quad}
\def\ou{\quad \textrm{or} \quad}
\def\boite{\hskip 0.5mm {\Box} \hskip 0.5mm}


\DeclareMathOperator\PPE{PPE}
\DeclareMathOperator\vNM{vNM}
\DeclareMathOperator\argmax{argmax}
\DeclareMathOperator\Prob{Prob}
\DeclareMathOperator\supp{supp}
\DeclareMathOperator\Eq{Eq}
\DeclareMathOperator\inte{int}
\DeclareMathOperator\co{co}
\DeclareMathOperator\cl{cl}
\DeclareMathOperator\Dirac{Dirac}
\DeclareMathOperator\F{F}
\DeclareMathOperator\As{As}
\DeclareMathOperator\Ir{Ir}
\DeclareMathOperator\MRS{MRS}
\DeclareMathOperator\Range{Im}
\DeclareMathOperator\DIV{DIV}
\DeclareMathOperator\A{D}
%------------------------------ theoren styles

\theoremstyle{definition}
\newtheorem*{assumption}{Assumption}
\newtheorem*{conjecture}{Conjecture}


\theoremstyle{example}
\newtheorem*{proposition}{Proposition}
\newtheorem*{remark}{Remark}


%-----------------------------------------

\title[Corporate Finance]{(Some theoretical aspects of) \\ Corporate Finance}
\author[V. F. Martins-da-Rocha]{V. Filipe Martins-da-Rocha}
\institute[FGV]{Escola de P\'os-gradua\c{c}\~ao em Economia \\~\\ Funda\c{c}\~ao Getulio Vargas}
\date[3º trimestre, 2012]{Part 1.e. Stock Market Equilibrium with moral hazard: About competitive price perceptions}
\subject{Economic Theory}


\begin{document}

\frame{\titlepage}


%------------------------------------------------

\begin{frame}\frametitle{Working paper}

\begin{thebibliography}{8}


  \beamertemplatearticlebibitems

    \bibitem{BisinGottardiRutta}
    Bisin, A., Gottardi. P. and Rutta, G.
    \newblock Equilibrium Corporate Finance.
    \newblock {\em NYU Working Paper} 2009

    \bibitem{BisinGottardiRutta}
    Bisin, A., Gottardi. P. and Rutta, G.
    \newblock Equilibrium Corporate Finance: Makowski meets Prescott and Townsend
    \newblock {\em EUI Working Paper} 2011

    \end{thebibliography}

\end{frame}

%------------------------------------------------

\begin{frame}\frametitle{Markets}

\begin{itemize}
\item One firm $J=\{j\}$ (or a continuum of identical firms)
\item One riskless financial asset $A=\{a_0\}$
    \begin{itemize}
    \item The firm can issue debt (or short-sell the financial asset $a_0$) but cannot invest: $\beta \geq 0$
    \item Investors cannot issue debt but can invest (or buy the firm's debt): $\theta^i \geq 0$
    \end{itemize}
\item Stock market
    \begin{itemize}
    \item Investors cannot short-sell the firm's shares: $\alpha^i \geq 0$
    \end{itemize}
\end{itemize}

\end{frame}



%------------------------------------------------

\begin{frame}\frametitle{Technology}

\begin{itemize}
\item The production technology of the firm is defined as follows:
\[
Y^j \equiv \{(y_0,(y_s)_{s\in S}) \in \R \times \R^S \ \colon \ y_0 \leq 0 \et y_s\leq f_s(-y_0,\alert{e}) \}
\]
where
\[
f_s :[0,\infty) \times \alert{E} \to [0,\infty)
\]
\item The investment $-y_0$ is denoted by $k$ representing units of commodity invested in capital at $t=0$
\item The manager's effort $e\in E$ is private information
\end{itemize}

\end{frame}



%------------------------------------------------

\begin{frame}\frametitle{Firm's choices}

\begin{itemize}
\item The firm chooses investment $k$ and debt level $\beta \geq 0$
\item We allow for bankruptcy: the dividend $\delta_s$ in state~$s$ is
\[
[f_s(k,e) - \beta]^+ = \max\{f_s(k,e) -\beta,0\}
\]
\item The debt issued by the firm is then risky and the payoff is
\[
\min\{1,f_s(k,e)/\beta\}
\]
\item The manager is hired among investors and is offered a compensation package for exerting effort
\end{itemize}

We could alternatively let the manager choose $k$ and $\beta$ but we will show that if $k$ and $\beta$ are chosen to maximize the value of the firm, then we will get strong unanimity

\end{frame}

%------------------------------------------------------------------------------------------------

\begin{frame}\frametitle{Compensation package}

A manager's compensation package consists of
\begin{itemize}
\item a payment~$x_0\geq 0$ at $t=0$ in units of the consumption good
\item a portfolio of
    \begin{itemize}
    \item $\alpha^m \geq 0$ shares of the firm
    \item $\theta^m$ units of bonds
    \end{itemize}
\item The disutility of effort for agent~$j$ is represented by a function $v^j : E \to \R_+$
\end{itemize}

\end{frame}

%------------------------------------------------------------------------------------------------

\begin{frame}\frametitle{Budget restrictions for the chosen manager}

Assume that $i$ has been chosen as a manager
\begin{itemize}
\item we make the assumption that the consumption contract is exclusive
\item the manager will choose effort $e$ to maximize his expected utility
\end{itemize}

\[
u^i(c^i_0) + \beta^i \sum_{s\in S}\nu^i_s u^i(c^i_s) - v^i(e)
\]
where
\[
c^i_0 = \omega^i_0 + x_0
\]
and for each~$s$
\[
c^i_s = \omega^i_s + [f_s(k,e)-\beta]^+ \alpha^m + \min\{1,f_s(k,e)/\beta\} \theta^m
\]
We let $\widetilde{e}^i(k,\beta,\alpha^m,\theta^m,x_0)$ be the set of optimal effort levels given the firm's choice $(k,\beta)$ and the compensation scheme $(\alpha^m,\theta^m,x_0)$

\end{frame}

%------------------------------------------------

\begin{frame}\frametitle{Budget restrictions for the other agents}

Each investor~$i$ takes as given
\begin{itemize}
\item the firm's capital budgeting and financial strategy $(k,\beta)$
\item the manager's effort choice $e\in E$
\item the security price $q$ and the stock (equity) price $E$
\item the cost $W$ at $t=0$ of hiring a manager
\end{itemize}
and chooses
\begin{itemize}
\item a consumption plan $c=(c_0,c_1) \in \R_+ \times \R_+^S$
\item an investment $\theta^i\geq 0$ and a share $\alpha^i \geq 0$ satisfying
\end{itemize}

\end{frame}

%------------------------------------------------

\begin{frame}\frametitle{Budget restrictions for the other agents}

\begin{itemize}
\item the budget restriction at $t=0$
\[
c^i_0 + q \theta^i + E \alpha^i \leq \omega^i_0 + [E+q\beta - k - W] \xi^i
\]
\item the budget restriction at $t=1$ for each state $s$
\[
c^i_s \leq \omega^i_s + [f_s(k,e)-\beta]^+ \alpha^i + \min\{1,f_s(k,e)/\beta\} \theta^i
\]
\end{itemize}
\begin{block}{Opportunity cost}
We denote by $B^i(q,E,k,\beta,e,W)$ the set of all actions $(c,\theta,\alpha)$ satisfying the budget restrictions and we let $U^i(q,E,k,\beta,e,W)$ be the associated utility, i.e.,
\[
U^i(q,E,k,\beta,e,W) \equiv \sup \{ U^i(c) \ \colon \ (c,\theta,\alpha) \in B^i(q,E,k,\beta,e,W)\}
\]
\end{block}
\end{frame}

%------------------------------------------------

\begin{frame}\frametitle{Competitive equilibrium given the firm's decisions}

Fix
\begin{itemize}
\item a firm's decision $(k,\beta)$,
\item a choice of manager~$j$,
\item a compensation scheme $(\alpha^m,\theta^m,x_0)$
\item an effort level $e$
\end{itemize}
then a family
\[
\{(k,\beta,j,\alpha^m,\theta^m,x_0,e),(q,E),(c^i,\theta^i,\alpha^i)_{i\neq j}\}
\]
is a competitive equilibrium if

\end{frame}

%------------------------------------------------

\begin{frame}\frametitle{Competitive equilibrium given the firm's decisions}

\begin{itemize}
\item for every $(i\neq j)$ the action $(c^i,\theta^i,\alpha^i)$ is optimal where the hiring cost $W$ is given by
\[
(1-\xi^j) W = x_0 + E(\alpha^m-\xi^j) + q\theta^m - \xi^j[q \beta - k]
\]
\item the manager~$j$ chooses the optimal effort level $e$ given the compensation scheme
\item financial markets clear, i.e.,
\[
\theta^m + \sum_{i\neq j} \theta^i = \beta \et \alpha^m + \sum_{i\neq j} \alpha^i = 1
\]
\item consumption markets clear, i.e.,
\[
k + \sum_{i\in I} c^i_0 = \omega_0 \et \sum_{i\in I} c^i_1 = \omega_1 + f(k,e)
\]
\end{itemize}

\end{frame}


%------------------------------------------------

\begin{frame}\frametitle{Notations}

Fix a consumption plan  $c^i=(c^i_0,c^i_1)$ satisfying
\[
c^i_0 >0 \et c^i_s >0, \ \forall s\in S
\]
Given a random variable $w=(w_s)_{s\in S} \in \R^S$, we let
\[
\E^i[w] = \sum_{s\in S} \MRS^i_{s,0}(c^i) w_s
\]
where
\[
\MRS^i_{s,0}(c^i) \equiv \beta^i \nu^i_s \frac{\partial u^i(c^i_s)}{\partial u^i(c^i_0)}
\]

\end{frame}

%------------------------------------------------

\begin{frame}\frametitle{Competitive equilibrium: necessary conditions}

\begin{proposition}
Assume there exists a competitive equilibrium
\[
\{(k,\beta,j,\alpha^m,\theta^m,x_0,e),(q,E),(c^i,\theta^i,\alpha^i)_{i\neq j}\} \quad \textrm{with} \quad \beta >0
\]
Then
\begin{eqnarray*}
q & = & \max_{i\neq j} \sum_{s\in S} \MRS^i_{s,0}(c^i) \min\{1,f_s(k,e)/\beta\} \\
& = & \max_{i\neq j} \E^i \left[ \min\{1,f_1(k,e)/\beta\} \right]
\end{eqnarray*}
and
\begin{eqnarray*}
E & = & \max_{i\neq j} \sum_{s\in S} \MRS^i_{s,0}(c^i) [f_s(k)-\beta]^+ \\
& = & \max_{i\neq j} \E^i  \left[  f_1(k)- \beta \1_S \right]^+
\end{eqnarray*}
\end{proposition}

\end{frame}


%------------------------------------------------

\begin{frame}\frametitle{Investors' price perception: definition}

Fix a competitive stock market equilibrium
\[
\{(k,\beta,j,\alpha^m,\theta^m,x_0,e),(q,E),(c^i,\theta^i,\alpha^i)_{i\neq j}\}
\]
\begin{itemize}
\item If the firm's choices were different, say
\[
(\tilde{k},\tilde{\beta},\tilde{j},\tilde{\alpha}^m,\tilde{\theta}^m,\tilde{x}_0)
\]
\item then its cash flow would be $[f_s(\tilde{k},\tilde{e}) - \tilde{\beta}]^+$ where
 \[
 \tilde{e} \in \tilde{e}^{\tilde{j}} (\tilde{k},\tilde{\beta},\tilde{\alpha}^m,\tilde{\theta}^m,\tilde{x}_0)
 \]
\item and the debt repayment should be $\min\{\tilde{\beta}, f_s(\tilde{k},\tilde{e})\}$
\end{itemize}

Consequently, the stock and bond prices should be different

\end{frame}

%------------------------------------------------

\begin{frame}\frametitle{Investors' price perception: Makowski criteriom}

We assume that each investor
    \begin{itemize}
    \item is aware that the stock and bond prices is a function of
    \[
    \tilde{\phi} \equiv (\tilde{k},\tilde{\beta},\tilde{j},\tilde{\alpha}^m,\tilde{\theta}^m,\tilde{x}_0)
    \]
    \item believes (conjectures) that
        \[
        \widetilde{E}(\tilde{\phi},\tilde{e}) = \max_{i \alert{\in ?}} \E^i \left[ [f_1(\tilde{k},\tilde{e})-\tilde{\beta} \1_S]^+ \right]
        \]
        and
        \[
        \widetilde{q}(\tilde{\phi},\tilde{e}) = \max_{i \alert{\in ?}} \E^i \left[ \min\{\1_S,f_1(\tilde{k},\tilde{e})/\tilde{\beta} \} \right]
        \]
    \item where $\tilde{e}$ is the optimal effort of manager $\tilde{j}$ given the new compensation scheme
     \[
 \tilde{e} \in \tilde{e}^{\tilde{j}} (\tilde{k},\tilde{\beta},\tilde{\alpha}^m,\tilde{\theta}^m,\tilde{x}_0)
 \]
    \end{itemize}

\end{frame}


%------------------------------------------------------------------------------------------------

\begin{frame}\frametitle{Maximizing firm's value}

We let $V^j(\tilde{\phi})$ be the firm's value when hiring agent~$j$ as a manager
\[
V^j(\tilde{\phi}) \equiv \tilde{E}(\tilde{\phi},\tilde{e}) + \tilde{q}(\tilde{\phi},\tilde{e}) \tilde{\beta} - \tilde{k} - \tilde{W}(\tilde{\phi})
\]
where $\tilde{W}(\tilde{\phi})$ is the corresponding hiring cost
\[
(1-\xi^j) \tilde{W}(\tilde{\phi}) = \tilde{x}_0 + \tilde{E}(\tilde{\phi},\tilde{e}) [\tilde{\alpha}^m-\xi^j] + \tilde{q}(\tilde{\phi},\tilde{e}) \tilde{\theta}^m - \xi^j[\tilde{q}(\tilde{\phi},\tilde{e}) \tilde{\beta} - \tilde{k}]
\]
$\tilde{e}$ is the optimal effort of manager $j$ given the new compensation scheme
     \[
 \tilde{e} \in \tilde{e}^j (\tilde{k},\tilde{\beta},\tilde{\alpha}^m,\tilde{\theta}^m,\tilde{x}_0)
 \]
and the participation constraint is satisfied\footnote{Pose $V^j(\tilde{\phi})=-\infty$ if it is not satisfied}
\[
U^j(\tilde{c}^j) - v^j(\tilde{e}) \geq U^j(q,E,k,\beta,e,W)
\]
where $\tilde{c}^j$ is the utility of the manager associated to the compensation scheme
\end{frame}

%------------------------------------------------

\begin{frame}\frametitle{Strong unanimity: result}

\begin{theorem}[Strong unanimity]
Fix a competitive stock market equilibrium
\[
\{(\phi,e),(q,E),(c^i,\theta^i,\alpha^i)_{i\neq j}\} \quad \textrm{where} \quad \phi = (k,\beta,j,\alpha^m,\theta^m,x_0)
\]
where the firm has chosen $\phi$ to maximize its value
\[
\phi \in \argmax_{\tilde{\phi}} V^j(\tilde{\phi})
\]
among all possible managers
\[
j\in \argmax_{\tilde{j} \in I} \sup_{\tilde{\phi}} V^{\tilde{j}} (\tilde{\phi})
\]
Then there is strong unanimity among investors concerning the firm's choice
\end{theorem}

\end{frame}


%------------------------------------------------

\begin{frame}\frametitle{Constrained Pareto efficiency: admissible allocations}

\begin{definition}
An allocation $(c^i)_{i\in I}$ of consumption plans is said \alert{admissible} if
\begin{itemize}
\item it is \alert{feasible}: there exists an investment level $k\geq 0$ such that
\[
k + \sum_{i\in I} c^i_0 \leq \omega_0 \et \sum_{i\in I} c^i_1 \leq f(k) +  \sum_{i\in I} \omega^i_1
\]
\item it is \alert{attainable}: there exists $\beta\geq 0$ and $(\theta^i,\alpha^i)$ for each~$i$,
    \[
    c^i_s = \omega^i_s + [f_s(k,e) - \beta]^+ \alpha^i + \min\{1,f_s(k,e)/\beta\} \theta^i
    \]
\item it is \alert{incentive compatible}: there exist an agent~$j$ such that
\[
e \in \argmax_{\tilde{e} \in E} u^j(c^j_0) + \beta^j \sum_{s\in S} \nu^j_s u^j(c^j(\tilde{e})) - v^j(\tilde{e})
\]
where $c^j_s(\tilde{e}) = \omega^j_s + [f_s(k,\tilde{e}) - \beta]^+ \alpha^i + \min\{1,f_s(k,\tilde{e})/\beta\} \theta^i$
\end{itemize}
\end{definition}

\end{frame}

%------------------------------------------------

\begin{frame}\frametitle{Constrained Pareto efficiency}

\begin{remark}
Observe that every competitive equilibrium allocation is admissible
\end{remark}

\begin{definition}
An admissible allocation $(c^i)_{i\in I}$ is constrained Pareto efficient if we cannot find another admissible allocation $(\widetilde{c}^i)_{i\in I}$ which is a Pareto improvement, i.e.,
\[
\forall i\in I, \quad U^i(\widetilde{c}^i) \geq U^i(c^i)
\]
with at least one strict inequality
\end{definition}

\begin{theorem}
Every competitive equilibrium allocation is constrained Pareto efficient
\end{theorem}

\end{frame}



\end{document}


%------------------------------------------------

\begin{frame}\frametitle{}

\begin{itemize}
\item
\item
\end{itemize}

\end{frame}
